\sscp{Kowiai}{11-03-02}
There have been differing views about the classification of \ili{Kowiai},
sometimes by the same scholars.
For example, \citet[272]{Blust1993} says,
“Koiwai appears to be a CMP language,
but it is not closely related to the \ili{Onin}-\ili{Sekar}-\ili{Uruangnirin} group,
and its wider connections within CMP remain obscure.” Yet \citet{BlustTrussel2020}
assign \ili{Kowiai} under SHWNG.

\citet[80]{WalkerR1982} and \citet[22]{WalkerWalker1991} note
that the king’s family in \ili{Kowiai} traces its origins to the \ili{Gorom}
region off of eastern Seram;
and that \ili{Geser}-\ili{Gorom} is the closest language to \ili{Kowiai},
based on lexical similarity
(\citealp[84, 93]{WalkerR1982}, \citealp[120--122]{WalkerWalker1991}).

Looking at historical phonology
and the lexicon, \citet{Zobel2014} also proposed that \ili{Kowiai}’s
“closest relative” is \ili{Geser}
and then \ili{Watubela}, and presents data to support his claim.
The first are some lexically restricted innovations shown in
\trf{11-21}, supplemented by data from our own sources.

\begin{table}
	\centering\tcp{\tsc{Eastern Islands} lexical innovations}{11-21}
	\st{5.5pt}\begin{tabular}{lllllll}\lsptoprule
local gl.	&	wood	&	walk	&	roof thatch	&	no, not	&	steal	&	wet\\	\midrule
\ili{Geser}	&	ruk	&	taŋgi	&	barem	&	tei	&	fakaleus	&	abotan\\
\ili{Gorom}	&	du	&	tagi	&	barem	&	tei	&	haleus	&	kiribotan\\
\ili{Watubela}	&	lunam	&		&		&	teiʔ	&		&	kabetan\\
\ili{Kowiai}	&	ruʔ	&	na-taŋ	&	barem	&	tei	&	na-ɸaraʔeus	&	abotan\\
		\lspbottomrule
	\end{tabular}
\end{table}

\newpage
Apart from these lexical innovations,
\ili{Kowiai} shares both of the main phonological innovations that
define \tsc{Eastern Islands}: lexically specific split of *j
> **s, **l and merger of *d/*z/*R > **r.

Despite its secure classification as an \tsc{Eastern Islands}
language, \ili{Kowiai} also shares a number of innovations with neighbouring
Austronesian languages of the Bomberai Peninsula.
These changes include *ŋ > \it{g,
ŋg} and *l > \it{r} which are both shared with \tsc{East Bomberai}
(\srf{11-05-02}), as well as *y > \it{l}
which is shared with \ili{Uruangnirin} (\tsc{Bomberai Tip}, \srf{11-05-01}).
Nearby \ili{Teor} also has *ŋ > \it{g} (\srf{11-04}).

Examples of *l > \it{r} have been given in \trf{11-15}.
Examples of *ŋ > \it{g,
ŋg} are given in \qf{11-02}.
Word finally *ŋ > \it{n}.
There is also one example of word-initial *ŋ > \it{n}; *\tbr{ŋ}ajan
> \it{\tbr{n}esa} ‘name’, and one example of
*ŋ = \it{ŋ}; *sa\tbr{ŋ}a > \it{ai sa\tbr{ŋ}} ‘forked branch’.
\\

\noindent{\wg\st{6pt}\begin{tabularx}
{\linewidth}{>{\hspace{-\tabcolsep}}l<{\hspace{-\tabcolsep}}
>{\hspace{-\tabcolsep}\enspace}ll
>{\raggedright\arraybackslash}X}
\lsptoprule
%\tbx{00-2}	&	PMP	&	\ili{Kowiai}	&	local gl.	\\	\midrule
%\endfirsthead
\tbx{11-02}	&	PMP	&	\ili{Kowiai}	&	local gl.	\\	\midrule
%\endhead
	&	*də\tbr{ŋ}əR	&	na-ɸano\tbr{ŋg}ar	&	hear	\\
	&	*la\tbr{ŋ}it	&	ra\tbr{ŋg}it	&	sky	\\
	&	*Ra\tbr{ŋ}a	&	ra\tbr{ŋg}a{\tl}ra\tbr{ŋg}a	&	spider conch	\\
	&	*la\tbr{ŋ}uy	&	na-la\tbr{g}ur	&	swim (irregular *l = \it{l}, irregular final \it{r})	\\
	&	*\tbr{ŋ}adas	&	\tbr{g}aran	&	neck (PMP `palate')	\\
	&								&	\tbr{g}aran utaʔ	&	throat	\\
	&								&	\tbr{g}aran sus	&	soft palate	\\
	&	*sa\tbr{ŋ}əlaR	&	na-sa\tbr{g}ela 	&	fry (irr. *l = \it{l})	\\
	&	*təli\tbr{ŋ}a	&	teri\tbr{g}a	&	ear	\\
	&	*ta\tbr{ŋ}is	&	na-ta\tbr{g}is	&	cry	\\	\midrule
	&	*duyu\tbr{ŋ}	&	rui\tbr{n}	&	dugong	\\
	&	*diŋdi\tbr{ŋ}	&	reri\tbr{n}	&	wall	\\
	&	*təri\tbr{ŋ}	&	tiri\tbr{n}	&	bamboo species (\ili{Malay} `bambu Jawa')	\\
	&	*qajə\tbr{ŋ}	&	aru\tbr{n}	&	charcoal	\\
\lspbottomrule
\end{tabularx}\mbox{}\\}

%\noindent{\wg\st{6pt}\tblr{>{\hspace{-\tabcolsep}}l<{\hspace{-\tabcolsep}}
%>{\hspace{-\tabcolsep}\enspace}
%lll}
%\lsptoprule
%%\tbx{00-2}	&	PMP	&	\ili{Kowiai}	&	local gl.	\\	\midrule
%%\endfirsthead
%\tbx{11-02}	&	PMP	&	\ili{Kowiai}	&	local gl.	\\	\midrule
%%\endhead
	%&	*də\tbr{ŋ}əR	&	na-ɸano\tbr{ŋg}ar	&	hear	\\
	%&	*la\tbr{ŋ}it	&	ra\tbr{ŋg}it	&	sky	\\
	%&	*Ra\tbr{ŋ}a	&	ra\tbr{ŋg}a{\tl}ra\tbr{ŋg}a	&	spider conch	\\
	%&	*la\tbr{ŋ}uy	&	na-la\tbr{g}ur	&	swim (irregular *l = \it{l}, irregular final \it{r})	\\
	%&	*\tbr{ŋ}adas	&	\tbr{g}aran	&	neck (PMP `palate')	\\
	%&								&	\tbr{g}aran utaʔ	&	throat	\\
	%&								&	\tbr{g}aran sus	&	soft palate	\\
	%&	*sa\tbr{ŋ}əlaR	&	na-sa\tbr{g}ela 	&	fry (irr. *l = \it{l})	\\
	%&	*təli\tbr{ŋ}a	&	teri\tbr{g}a	&	ear	\\
	%&	*ta\tbr{ŋ}is	&	na-ta\tbr{g}is	&	cry	\\	\midrule
	%&	*duyu\tbr{ŋ}	&	rui\tbr{n}	&	dugong	\\
	%&	*diŋdi\tbr{ŋ}	&	reri\tbr{n}	&	wall	\\
	%&	*təri\tbr{ŋ}	&	tiri\tbr{n}	&	bamboo species (\ili{Malay} `bambu Jawa')	\\
	%&	*qajə\tbr{ŋ}	&	aru\tbr{n}	&	charcoal	\\
%\lspbottomrule
%\end{tabular}\mbox{}\\}

Examples of *y > \it{l} are given in \qf{11-03}.
Note that several of these examples involve prothesis of **y
to vowel initial words (see \srf{13-02-02} for more details).

\noindent{\wg\st{6pt}\begin{tabularx}
{\linewidth}{>{\hspace{-\tabcolsep}}l<{\hspace{-\tabcolsep}}
>{\hspace{-\tabcolsep}\enspace}llll
>{\raggedright\arraybackslash}X}
\lsptoprule
%\tbx{00-3}	&	PMP	&	Intermediate forms	&	\ili{Kowiai}	&	local gl.	\\	\midrule
%\endfirsthead
\tbx{11-03}	&	PMP	&	\mc{2}{l}{Intermediate forms}	&	\ili{Kowiai}	&	local gl.	\\	\midrule
%\endhead
	&	*wahi\tbr{y}əR	&		&	**wa\tbr{y}ər	&	wa\tbr{l}ar	&	water	\\
	&	*ma-hə\tbr{y}aq	&		&	**ma\tbr{y}a	&	na-ma{\tl}ma\tbr{l}a	&	shy, ashamed	\\
	&	*ba\tbr{y}awak	&	**ba\tbr{y}ak	&	**\tbr{y}abak	&	\tbr{l}aɸaʔos	&	monitor lizard (C met.)	\\
	&	*i-aku	&		&	**\tbr{y}aku	&	\tbr{l}aʔ	&	\tsc{1sg}	\\
	&	*añam	&		&	**\tbr{y}añam	&	na-\tbr{l}alan	&	weave (baskets, mats) (perhaps *ñ > \it{l}?)	\\
	&	*qatay	&	**ata	&	**\tbr{y}ata	&	\tbr{l}ata	&	liver	\\
	&	*qapuR	&	**apuR	&	**\tbr{y}apur	&	\tbr{l}aɸor	&	calcium	\\
	&	*qazay	&	**azay	&	**\tbr{y}ara	&	\tbr{l}araʔai	&	jaw	\\
	&	*hapuy	&	**apuy	&	**\tbr{y}api	&	\tbr{l}aɸ	&	fire	\\
\lspbottomrule
\end{tabularx}\mbox{}\\}
